%% kapa: exam=1 round=27 year=2016 period=2 form=A subject=accounting subject_ko=회계학 seq=21 number=61
%% source: Q-Net 감정평가사 기출문제 (cst00309) artlSeq=5169402 file=extracted/1008/5169402/제27회 감정평가사 1차 문제/제27회 감정평가사 2교시문제(A형).pdf
%% conversion: kordoc 4.24.51 -> tools/gen1.py (auto); review=proofread-sonnet; flags=none
\begin{question}{61}
(주)감평은 20×6년 1월 1일 본사 건물을 새로 마련하기 위한 공사계약을 체결하였다. 본 건물은 적격자산에 해당되며 차입원가 자본화와 관련된 내용은 다음과 같다.
\begin{framedtext}
○ 공사 착공일: 20×6년 5월 1일 \par
○ 공사 완공일: 20×6년 12월 31일 \par
○ 공사대금 지급 \par
\hspace*{0.5em}- 20×6년 5월 1일 ￦1,200,000 \par
\hspace*{0.5em}- 20×6년 9월 1일 ￦900,000 \par
\hspace*{0.5em}- 20×6년 10월 1일 ￦500,000 \par
○ 차입금 \par
\hspace*{0.5em}- 20×6년 2월 1일 본사 건물 공사와 관련하여 상환일이 20×7년 12월 31일인 ￦1,200,000(연 8\%)을 차입 \par
\hspace*{0.5em}- 20×6년 7월 1일 건물 공사 이외의 목적으로 상환일이 20×7년 6월 30일인 ￦800,000(연 10\%)을 차입 \par
\hspace*{0.5em}- 20×6년 10월 1일 건물 공사 이외의 목적으로 상환일이 20×8년 12월 31일인 ￦1,500,000(연 12\%)을 차입 \par
\hspace*{0.5em}- 20×6년 2월 1일 차입한 금액 중에서 ￦120,000을 20×6년 4월 1일부터 8월 31일까지 연 5\% 수익률로 일시적 투자 \par
\end{framedtext}
20×6년 12월 31일 재무상태표에 계상되는 건물의 취득원가는? (단, 계산시 월할계산하며 단수차이가 있으면 가장 근사치를 선택한다.)
\begin{choices}
  \item ￦2,671,214
  \item ￦2,685,613
  \item ￦2,697,000
  \item ￦2,702,123
  \item ￦2,713,000
\end{choices}
\end{question}
\answer{5}
